
\lesson{02/11}
\subsection{Characters}
\textbf{Def.} The \textit{character} of a representation $U_\alpha$ of a finite
group $G$ is the map
\begin{align*}
    \chi_\alpha:&G\to\mathbb{C}\\
                &g\to\chi_\alpha(g):=\tr(U_\alpha(g))
\end{align*}
\textbf{Observation 1} If two representations are linked by a similarity
transfomation
$$D_\alpha(g)=AD_\beta(g)A^{-1}$$
their characters are the same, for the properties of the trace
$$\chi_\alpha(g)=\chi_\beta(g)$$
\textbf{Observation 2} If $g_1,g_2\in G$ are so that
$$\exists g\in G\text{ s.t. }g_2=g\cdot g_1\cdot g^{-1}$$
then
$$\chi_\alpha(g_1)=\chi_\alpha(g_2),\,\forall \alpha$$
and $g_1,g_2$ are in the same \textit{conjugacy class}, which is an equivalence
class.

\textbf{Observation 3} The character of any representation for any group
evaluated at the identity yields
$$\chi_\alpha(e)=\dim(D_\alpha),\,\forall G,\,\forall D_\alpha$$
\textbf{Observation 4} From the orthogonality theorem follows that
$$\sum_{g\in G}^{} \chi_\alpha^*(g)\chi_\beta(g) = |G|\delta_{\alpha\beta}$$
\textbf{Observation 5} For a completely reducible representation, i.e.\footnote{
unclear, maybe it is explained in lessons 6,7,8}
$$U=\bigoplus_\alpha n_\alpha U_\alpha$$
the linearity of the trace implies
$$\tr U(g)=\sum_{\beta}^{}n_\beta\tr(U_\beta(g))\implies\chi(g)=
\sum_{\beta}^{}n_\beta\chi_\beta(g)$$
multiplying by $\chi^*_a(g)$ and summing over $g\in G$ yields
$$\sum_{g\in G}^{}\chi_\alpha^*(g)\chi(g) = \sum_{\beta}^{}n_\beta\sum_{g\in G}
\chi_\alpha^*(g)\chi_\beta(g)=\sum_{\beta}^{}n_\beta|G|\delta_{\alpha\beta}=
n_\alpha|G|$$
thus
$$n_\alpha = \frac{1}{|G|}\sum_{g\in G}^{}\chi_\alpha^*\chi(g)$$

By considering the conjugacy equivalence classes $c_i$, $i=1,2,...k$,
this becomes
$$n_\alpha= \frac{1}{|G|}\sum_{i=1}^{k}|c_i|\chi_\alpha^*(c_i)\chi(c_i).$$\\

By defining the \textit{character vector}
$$\vec{\chi}_\alpha = \begin{pmatrix}
\sqrt{|c_1|}\chi_\alpha(c_1)\\...\\\sqrt{|c_k|}\chi_\alpha(c_k)
\end{pmatrix}\in\mathbb{C}^k
$$
one can interpret the orthogonality condition literally: this implies that
for any finite group \textbf{the number of irreps is at most (actually is exactly)
the number of conjugacy classes.}

\textbf{Theorem} \textit{If $G$ is Abelian and finite, all its irreps have
dimension 1}\\
\textbf{Proof} If $G$ is Abelian, then all conjugacy classes must contain a single
element:
$$g_1\sim g_2 \iff g_1 = g\cdot g_2\cdot g^{-1}=g\cdot g^{-1}\cdot g_2=g_2$$
So $|c_i|=1$ and
$$|G|=\sum_{i=1}^{k}|c_i| = k $$
so the number of irreps is $|G|$, but, by the Burnside theorem
$$|G|=\sum_{\alpha}^{}(\dim(V_\alpha))^2$$
since $|G|=\text{ number of irreps }=\sum_{\alpha}^{}1$,
$$\sum_{\alpha}^{}(\dim(V_\alpha))=\sum_{\alpha}^{}1\iff\dim(V_\alpha)=1,\,
\forall\alpha$$\qed

\textbf{Note} This also holds for infinite groups.

\textbf{Example: $S_3$}

By Burnside's theorem, the irreps of $S_3$ can have dimensions
$$1,1,1,1,1,1\text{ or }1,1,2.$$
In $S_3$ permutations are of three kinds:
\begin{enumerate}
    \item One permutation with no non-trivial cycles;
    \item Three permutations that are cycles of length 2;
    \item Two permutations that change every element.
\end{enumerate}
It can be proven that these correspond to the conjugacy classes, so there exist
only $3$ irreps, which, by Burnside, have dimensions $1,1,2$.

The characters of these $3$ representations are
\begin{table}[h]
    \centering
    \begin{tabular}{c|c|c|c}
            &$\chi(E)$&$\chi(1\to2)$&$\chi(1\to2\to3)$\\ \hline
        $1$ &   1&1&1\\\hline
        $1'$&   1&$a$&$b$\\\hline
        $2$ &   1&$c$&$d$
\end{tabular}
\end{table}\newpage
$1$ is the trivial $1-$dimension representation ($\forall g\in G, U(g)=1\in
\mathbb{R}$), while $1'$ maps all elements to their parity (in the sense of
permutations).

Some properties of $a,b,c,d$ can be deduced from the character vectors:
$$\vec{\chi}_1=\begin{pmatrix}
1\\\sqrt{3}\\\sqrt{2}
\end{pmatrix},\; \vec{\chi}_{1'}=\begin{pmatrix}
1\\\sqrt{3}a\\\sqrt{2}b
\end{pmatrix},\; \vec{\chi}_{2}=\begin{pmatrix}
1\\\sqrt{3}c\\\sqrt{2}d
\end{pmatrix}
$$
since they are orthogonal, it follows that
$$
\begin{dcases*}
\vec{\chi}_1^\dagger\cdot\vec{\chi}_{1'} = 1+3a+2b=0\\
\vec{\chi}_1^\dagger\cdot\vec{\chi}_2 = 2+3c+2d=0\\
\vec{\chi}_{1'}^\dagger\cdot\vec{\chi}_2 = 1+3ac+2bd=0\\
\vec{\chi}_1\cdot\vec{\chi}_1\equiv 6\\
\vec{\chi}_{1'}\cdot\vec{\chi}_{1'} = 1+3a^2+2b^2 = 6\\
\vec{\chi}_2\cdot\vec{\chi}_2 = 4+3c^2+2d^2 = 6
\end{dcases*}
$$
Since the $1'$ representation maps the elements to their parity, $a=-1$ and $b=1$,
so the system becomes
$$\begin{dcases*}
1-3+2=0\\
2+3c+2d=0\\
1-c+2d=0\\
6=6\\
4+c^2+2d^2=6
\end{dcases*}\implies
\begin{dcases*}
3c+2d = -2\\
-3c+2d = -2\\
c^2+2d^2 = 2
\end{dcases*}\implies\begin{dcases*}
d = -1\\
c = 0\\
0^2+2(-1)^2 = 2
\end{dcases*}
$$
\begin{table}[h]
    \centering
    \begin{tabular}{c|c|c|c}
            &$\chi(E)$&$\chi(1\to2)$&$\chi(1\to2\to3)$\\ \hline
        $1$ &   $1$&$1$&$1$\\\hline
        $1'$&   $1$&$-1$&$1$\\\hline
        $2$ &   $1$&$0$&$-1$
\end{tabular}
\end{table}\\

Consider now a representation of $S_3$:
$$U(e) = \begin{pmatrix}
    1&0&0\\0&1&0\\0&0&1
\end{pmatrix},\,U(1\to2)=\begin{pmatrix}
    0&1&0\\
    1&0&0\\
    0&0&1
\end{pmatrix},\,U(1\to3)=\begin{pmatrix}
    0&0&1\\
    0&1&0\\
    1&0&0
\end{pmatrix},\,U(2\to3)=\begin{pmatrix}
    1&0&0\\
    0&0&1\\
    0&1&0
\end{pmatrix}
$$
$$U(1\to2\to3) = \begin{pmatrix}
    0&0&1\\
    1&0&0\\
    0&1&0
\end{pmatrix},\,U(1\to3\to2) = \begin{pmatrix}
    0&1&0\\
    0&0&1\\
    1&0&0
\end{pmatrix}
$$
Since all these matrices are unitary and the representation is reducible, then
it is completely reducible. The character vector of $U$ is
$$\vec{\chi}_U=\begin{pmatrix}
1\cdot 3\\\sqrt{3}\cdot1\\0
\end{pmatrix}
$$
For the previous results,
\begin{gather*}
n_1 = \frac{1}{6}\sum_{i=1}^{3}|c_i|\chi_1^*(c_i)\chi_U(c_i)=
\frac{1}{6}\vec{\chi}_1^\dagger\cdot\vec{\chi}_U = [...]=1\\
n_{1'}= \frac{1}{6}\vec{\chi}_{1'}^\dagger\cdot\vec{\chi}_U = [...]=0\\
n_2 = \frac{1}{6}\vec{\chi}_2^\dagger\cdot\vec{\chi}_U = [...] = 1
\end{gather*}
So this means that the representation is reducible to
$$U = 1\oplus2$$\\

\textbf{Def.} A \textit{left-regular representation} of a finite group
$G$ is a representation which acts on $G$ as a left translation. More
precisely, consider the vector space $V=\mathbb{C}^{|G|}$ with base
$\left\{e_g\right\}_{g\in G}$: the left translation is defined as
\begin{align*}
    L:&G\to\text{Perm}(G)\\
      &g\to L_g\text{ s.t. }L_g(h)=g\cdot h
\end{align*}
and the regular representation is
\begin{align*}
&D(g):V\to V\\
&e_h\to D(g)e_h=e_{g\cdot h}\\
&D\left(g\right)v = D(g)\left(\sum_{h\in G}^{}v^he_h\right):=\sum_{h\in G}
v^hD(g)e_h = \sum_{h\in G}^{}v^h e_{g\cdot h}
\end{align*}
In matrix form the element corresponding to $e_h,e_{h'}$ is
$$D(g)_{h,h'}=\begin{cases}
    1&\text{if }h=g\cdot h'\\
    0&\text{otherwise}
\end{cases}
$$
this means that for $g=e$, the matrix is the identity of $V$; on the other hand,
if $g\neq e$, the diagonal is always $0$, which means all $D(g\neq e)$ are 
\textit{traceless}:
$$\forall g\in G,\, \chi(g)=|G|\delta_{ge}.$$\\

\textbf{Def.} A \textit{right-regular representation} of a finite group $G$ is a
representation which acts on $G$ as a right translation:
\begin{align}
&R:G\to\text{Perm}(G)\\
& g\to R_g\text{ s.t. } R_g(h)=h\cdot g^{-1}
\end{align}
and the representation is similarly defined.

\section{Tensors}
\textbf{Def.} Given a vector space $V$ of dimension $n$ over the complex field,
the \textit{dual vector space} $V^*$ is
$$V^*:=\left\{f:V\to\mathbb{C}\text{ s.t. $f$ is linear}\right\}$$
$V^*$ is a vector space with $\dim(V^*)=\dim(V)$.

Given a basis $\left\{e_i\right\}\subset V$, a natural choice of basis for 
$V^*$ is the \textit{dual basis}
$$\left\{e^{*i}\right\}\text{, s.t. }e^{*i}(e_j)=\delta\indices{^i_j}$$

If $V$ is equipped with a scalar product $\langle{\cdot},{\cdot}\rangle$ and the
basis is orthonormal, this dual basis is precisely
$$e^{*i}:=\langle e_i,\cdot\rangle\to e^{*i}(e_j):=\langle e_i,e_j\rangle =
\delta\indices{^i_j}$$
which is actually where the notations $e^{i*}=\bra{e_i}$ and $e_j=\ket{e_j}$
come from.
